\chapter[CONCLUSÃO]{CONCLUSÃO}
\label{chap:conclusao}

Este trabalho apresentou o desenvolvimento e a avaliação de um sistema automatizado de apoio à detecção de melanoma, integrando máscaras de segmentação de referência, modelos de linguagem multimodal e uma arquitetura multiagente coordenada via LangGraph. Mais do que a comparação de acurácia com um \textit{baseline} supervisionado, a principal contribuição metodológica deste trabalho reside em duas frentes: um \textbf{framework de cinco eixos} para avaliação automatizada da consistência estrutural e lógica de relatórios diagnósticos gerados por LLMs, e uma \textbf{análise de calibração de incerteza} que revelou que o modelo de melhor desempenho concentra sua hesitação diagnóstica justamente nos casos clinicamente mais ambíguos. A abordagem proposta visa complementar o raciocínio clínico humano, oferecendo diagnósticos estruturados, explicáveis e auditáveis a partir de imagens dermoscópicas do conjunto PH2 \cite{mendonca2013ph2}.

A adoção das máscaras de segmentação de referência (\textit{ground truth}) do PH2 como entrada ao sistema, em lugar de segmentações automáticas, foi uma decisão metodológica central que aumentou a validade interna dos experimentos. Ao eliminar a variabilidade introduzida por modelos de segmentação, foi possível atribuir as diferenças de desempenho exclusivamente à capacidade de raciocínio das LLMs, isolando o componente de interesse para avaliação. O conjunto PH2, composto por 200 imagens dermoscópicas (40 melanomas e 160 não-melanomas) com anotações especializadas, mostrou-se adequado para a prova de conceito proposta.

A aplicação do framework de cinco eixos revelou o achado mais informativo desta pesquisa. Ambos os modelos atingiram 100\% de completude estrutural, mas o Gemini-3.1-flash-lite demonstrou desempenho superior nos demais eixos: 100\% de acerto aritmético no protocolo ABCD (vs.\ 96\% do GPT), 46\% de aderência plena ao limiar declarado (vs.\ 6\%), 78\% de consenso inter-protocolos (vs.\ 46\%) e 85,2\% de acurácia nos casos com diagnóstico definitivo (vs.\ 22,2\%). De forma decisiva, o Gemini-3.1-flash-lite exibiu \textbf{calibração de incerteza clinicamente coerente}: seus diagnósticos ``suspeitos'' e ``indeterminados'' concentraram-se justamente nos casos com maior prevalência real de melanoma, evidenciando que o modelo hesita nos casos genuinamente difíceis. O GPT-5.4-nano, por sua vez, emitiu diagnóstico não-definitivo em 82\% dos casos e apresentou calibração invertida, os casos classificados como indeterminados correspondiam a 56,1\% de melanomas reais, indicando que o modelo não discrimina adequadamente os casos de maior risco. Esse contraste demonstra que o framework de cinco eixos captura diferenças de confiabilidade clínica entre modelos que a acurácia binária, isoladamente, não revela.

Quanto à detecção binária de melanoma, o \textit{baseline} supervisionado (ResNet-50), treinado diretamente sobre o PH2 com divisão estratificada 70/30 e ponderação de classes na função de perda, obteve acurácia de 0,90, precisão de 0,57, sensibilidade de 1,00 (identificando corretamente todos os melanomas do conjunto de teste) e F1-score de 0,73 sobre seu conjunto de teste (60 imagens, das quais 8 de melanoma), superando os modelos de linguagem multimodal em sensibilidade, resultado esperado para modelos multimodais generalistas sem ajuste fino especializado em dermatoscopia. O GPT-5.4-nano \cite{openai2025gpt5nano} atingiu acurácia de 0,68 e F1-score de 0,68, enquanto o Gemini-3.1-flash-lite \cite{google2025gemini31} obteve acurácia de 0,76 e F1-score de 0,74, ambos avaliados sobre um subconjunto balanceado e distinto do PH2 (50 imagens, 25 melanoma e 25 não-melanoma); como acurácia, precisão e F1-score dependem da prevalência de classes, esses valores não são diretamente comparáveis entre os dois conjuntos de teste, cuja diferença de composição é discutida como limitação no Capítulo~5. Ainda assim, as LLMs introduzem uma dimensão interpretativa ausente em classificadores tradicionais: a capacidade de expressar explicitamente os critérios diagnósticos utilizados na decisão, produzindo relatórios auditáveis e rastreáveis, dimensão que o framework de cinco eixos permite justamente mensurar de forma objetiva.

A arquitetura multiagente, ao consolidar pareceres independentes de três agentes especializados em protocolos clínicos distintos (ABCD, Menzies e SPCL) em um relatório unificado pelo agente sumarizador, demonstrou utilidade qualitativa relevante: mesmo quando o diagnóstico final era incorreto, as justificativas geradas eram semanticamente coerentes e rastreáveis. Esse comportamento aproxima o sistema de um processo de segunda opinião médica distribuída, aumentando a robustez interpretativa sem substituir a necessidade de modelos visuais especializados.

Os resultados confirmam que o sistema proposto funciona como uma \emph{prova de conceito} válida, demonstrando a viabilidade técnica de combinar imagens dermoscópicas, segmentação por referência, LLMs multimodais e estruturas multiagentes para a detecção de melanoma com diagnósticos explicáveis. Mais relevante do que o desempenho classificatório bruto, cuja comparação direta entre a ResNet-50 e as LLMs é limitada pela diferença de prevalência entre os conjuntos de teste utilizados, é a demonstração de que o framework de cinco eixos e a análise de calibração de incerteza constituem um protocolo reprodutível para auditar a confiabilidade do raciocínio diagnóstico de LLMs, revelando diferenças de qualidade entre modelos que passariam despercebidas em uma avaliação baseada exclusivamente em métricas de classificação. Ainda assim, o trabalho expõe limites claros que motivam investigações futuras, em especial quanto à acurácia classificatória das LLMs e à escala dos experimentos.

Como direções para trabalhos futuros, propõe-se:
\begin{itemize}
    \item \textit{Fine-tuning} supervisionado dos modelos LLM em bases dermoscópicas especializadas, visando melhorar o desempenho na detecção de melanoma;
    \item Integração de um módulo de segmentação automático ao \textit{pipeline}, substituindo as máscaras GT por predições de modelos U-Net ou similares, para avaliação do impacto da qualidade da segmentação no diagnóstico final;
    \item Expansão para conjuntos de dados maiores e mais diversificados (e.g., ISIC, HAM10000), permitindo análises estatísticas mais robustas e maior generalização;
    \item Estimativa de desvio-padrão e intervalos de confiança via \textit{bootstrap}, condicionada à ampliação do número de imagens de melanoma disponíveis para teste;
    \item Estudo de ablação das estratégias de \textit{prompting} (Chain-of-Thought, Risk Assessment e Consistency Prompt), atualmente combinadas em um único \textit{prompt} composto, para mensurar a contribuição individual de cada uma sobre o desempenho diagnóstico e a consistência textual;
    \item Ampliação do escopo diagnóstico para além do melanoma, incluindo outras neoplasias cutâneas como carcinoma basocelular e carcinoma espinocelular;
    \item Avaliação clínica prospectiva do sistema por profissionais de dermatologia, visando à validação em cenários reais de triagem.
\end{itemize}

A integração de LLMs multimodais e sistemas multiagentes para a detecção de melanoma representa uma abordagem ainda incipiente na literatura, com potencial significativo para democratizar o acesso ao diagnóstico assistido por IA em regiões com escassez de especialistas. O sistema proposto, apesar de suas limitações atuais, estabelece uma base metodológica sólida e replicável para investigações futuras nessa interseção entre inteligência artificial e medicina clínica.
